\chapter{Método do Trabalho}\label{cap:metodo}
% ===========================================================================

Este capítulo descreve o processo adotado para desenvolver o projeto, cujo
objetivo é implementar um protótipo de livraria digital com preservação de
privacidade por meio de Oblivious Transfer (OT). O método foi estruturado em
fases de forma a manter alinhamento com os objetivos definidos no \autoref{cap:introducao}.

\section{Visão Geral do Processo}

O processo de desenvolvimento foi organizado em seis fases principais:
fundamentação teórica, especificação de requisitos, projeto da solução,
implementação do protótipo, testes e avaliação, e consolidação dos resultados.

\section{Fase 1 --- Fundamentação Teórica}

Na primeira fase, foram estudados os conceitos criptográficos necessários para
o projeto, com ênfase em OT e suas variantes (como 1-out-of-n e GOT), além de
técnicas relacionadas, como OT Extension e fundamentos do criptossistema de
Rabin. Essa etapa forneceu base para as decisões arquiteturais e para a
definição dos limites de segurança do protótipo \cite{rabin2005, tassa2011}.

\section{Fase 2 --- Especificação de Requisitos}

Com base nos objetivos do trabalho, foi conduzida a elicitação e organização
dos requisitos funcionais e não funcionais do sistema. Essa fase formaliza o
que a solução deve fazer e quais restrições de desempenho, segurança e
usabilidade devem ser atendidas. O detalhamento dessa etapa é apresentado no
\autoref{cap:especificacao}.

\section{Fase 3 --- Projeto da Solução}

Após a definição dos requisitos, pretende-se fazer o projeto com uma
arquitetura inicial de referência. Os objetivos iniciais do projeto são:
organização em três componentes (sistema web, CDN e programa local), modelagem
preliminar do fluxo de compra privada, alternativas de agrupamento de itens
por faixa de preço e possíveis formatos de artefatos de criptografia enviados.
Nessa etapa, os critérios de separação entre as operações do sistema web e as
operações locais do cliente serão definidos conforme os resultados parciais. O
detalhamento dessa etapa é apresentado no \autoref{cap:desenvolvimento}.

\section{Fase 4 --- Implementação Incremental}

A implementação deve ser conduzida em incrementos. A sequência abaixo
representa uma proposta inicial, passível de ajustes durante o
desenvolvimento:

\begin{itemize}
\item estrutura de catálogo e cadastro de itens;
\item fluxo de compra com coleta mínima de dados do cliente;
\item execução do mecanismo criptográfico associado à escolha privada;
\item emissão de tokens com assinatura cega e entrega do conteúdo cifrado pela
CDN;
\item suporte ao processo de recuperação local dos itens adquiridos.
\end{itemize}

Esse encadeamento foi definido como estratégia de partida e poderá ser
reordenado caso surjam problemas técnicos ou necessidades de validação não
previstas. O detalhamento dessa etapa é apresentado no \autoref{cap:desenvolvimento}.

\section{Fase 5 --- Testes e Avaliação}

Na fase de avaliação, pretende-se priorizar três questões: (i) eficácia do
sistema, com foco no atendimento aos requisitos; (ii) nível de privacidade,
buscando provas de que o servidor não infere a escolha específica do cliente;
e (iii) desempenho, com medições de tempo de busca, compra e entrega. Os
procedimentos exatos de teste e o recorte final de métricas serão definidos ao
longo da implementação e posteriormente descritos no \autoref{cap:desenvolvimento}.

\section{Fase 6 --- Consolidação e Análise Final}

Na etapa final, pretende-se consolidar os resultados alcançados e analisá-los
com os objetivos do projeto, dando ênfase nas limitações observadas,
contribuições efetivas e possibilidades de continuidade. O nível de
profundidade dessa análise dependerá do que for possível concluir nas etapas
anteriores e será apresentado no capítulo de considerações finais.
